\scanpage{154}\mspage{9}
3) On peut faire les ${=}\mathrm{s}$ choses pour des coeff tels que $\mathbb{Z}_{\ell}$ (Jouanolou, thèse -- non publiée je crains !)

\inserted{4) Ça marche pour $f:X\unsure{\dashv|\longrightarrow}S$ \unsure{prés.|pres.} finie séparé, si $f$ a des propriétés de ``locale trivialité cohomologique'' (propr. locales \uline{en haut}) p.ex. $f$ \uline{lisse} ; on trouve $\underline{H}_{*X/S}=Rf_{!}\,f^{!}(\Lambda_S)$.}

B) Artin et Mazur ont étudié, dans un esprit voisin du tien, \uline{l'autodualité} de la jacobienne d'une courbe relative $X/S$. Il faut les consulter pour les résultats précis, peut-être cela marche dès que $X/S$ est propre et plat de dim relative $1$ -- en tous cas c'est OK sur un anneau de valuation discrète à fibre \uline{générique} lisse. La fibre spéciale peut être très sauvage. (J'ai utilisé \deleted{\unsure{pour|leur}}leur résultat dans SGA \deleted{\unsure{3}}7 pour prouver, dans le cas des jacobiennes, une conjecture de dualité sur les groupes de composantes connexes associés \unsure{aux} \deleted{\unsure{de|des} \unsure{V.A.|VA.}} modèles de Néron de V.A. duales l'une de l'autre....)

\figure{\begin{tikzpicture}
\draw (0,0)--(1.80em,-0.025em);
\draw (.63em,-.29em)--(1.63em,-.29em);
\end{tikzpicture}}

\deleted{\unsure{C)}} Vers la fin des années 50\inserted{(ou début des 60 ?)}, alors que les grands fourbis cohomologiques ($f_{!},f^{!}$, coh. étale etc) sortaient à peine des limbes, le cours de Serre sur la théorie de \unsure{Rosenlicht} et \unsure{Lang} des \unsure{Jacobiennes|jacobiennes|Jacobienne}
% archive sheet number: 154
